Figure~\ref{fig:atlas-oscillating-sequence} shows why two subsequences with different limits prevent convergence.

\begin{figure}[!htbp]
\centering
\begin{tikzpicture}[>=Stealth]
\begin{axis}[width=11cm,height=5.8cm,axis lines=middle,ticklabel style={font=\small},label style={font=\small},legend style={font=\scriptsize,draw=black!20,fill=white},xlabel={$n$},ylabel={$a_n$},xmin=0,xmax=13,ymin=-1.5,ymax=1.5,xtick={2,4,6,8,10,12},ytick={-1,1}]
\addplot[black!30,dashed,domain=0:12.5] {1};
\addplot[black!30,dashed,domain=0:12.5] {-1};
\addplot[only marks,figblue,mark=*,mark size=2pt] coordinates {(2,1)(4,1)(6,1)(8,1)(10,1)(12,1)};
\addplot[only marks,figred,mark=square*,mark size=2pt] coordinates {(1,-1)(3,-1)(5,-1)(7,-1)(9,-1)(11,-1)};
\end{axis}
\end{tikzpicture}
\caption{The sequence $a_n=(-1)^n$ alternates between two horizontal levels. The even subsequence is constantly $1$ and the odd subsequence is constantly $-1$, so the full sequence cannot approach a single limit.}
\label{fig:atlas-oscillating-sequence}
\end{figure}